% Written by Petr Gotthard
% Codepage ISO-8859-2

\section{Matematika a logika}
\begin{multicols}{2}

\noindent
Generální ředitel jednoho velkého podniku měl letět odpoledne
do vzdáleného města, kde byl pobočný závod. Ráno však přišel
k~němu noční vrátný, známý svou pověrčivostí, a začal
ho z~této cesty zrazovat. ,,Dnes se mně zdálo, že se letadlo,
kterým chcete letět, zřítilo,'' říkal. Generální ředitel vrátného
odbyl, ale pak mu to přece nedalo, myslel stále na padající letadlo
a nakonec cestu odložil. Druhého dne se v~novinách dočetl,
že se letadlo, kterým měl letět, skutečně zřítilo. Zavolal vrátného
a v~pohnutí mu děkoval za záchranu života. ,,Dávám vám
jako odměnu 5000 Kč, ale protože mám pevné zásady, zároveň vás
propouštím.'' Proč dostal vrátný výpověď?\\[1 mm]
{\sl Noční vrátný řekl ráno, že se mu dnes zdálo o letadlu, spal
tedy ve službě, a proto byl propuštěn.}\\

\noindent
Hotel Continental v~Chicagu byl proslulý hotelovými
lupiči. Jeden lupič byl přistižen takto: V~pokoji číslo
54 seděl host Max Brown a četl The New York Times. Najednou se
ozvalo slabší zaklepání, a něž se Max Brown ozval, otevřely se
dveře a v~nich se objevil navoněný elegán. Když spatřil
sedícího muže u stolu, zrozpačitěl a pak se narychlo omlouval:
,,Promiňte pane, spletl jsem si dveře svého pokoje.`` Max Brown
nedal na tuto výmluvu, vyvolal rychle poplach, a tak byl lupič
zadržen. Podle čeho Max Brown poznal, že se elegán nespletl ve
dveřích?\\[1 mm]
{\sl Max Brown poznal, že nejde o hosta, který si spletl dveře
pokoje, podle toho, že elegán před vstupem klepal, čímž chtěl
zřejmě zjistit, je-li někdo uvnitř. Při vstupu do vlastního pokoje
by nikdy neklepal.}\\

\noindent
Měsíc v~první čtvrti byl právě na obzoru, když
se automobil bez rozsvícených světel prořítil vysokou rychlostí
vesnicí. Byl však zastaven policií a řidič musel zaplatit pokutu.
Byl důvod k~pokutě jeden nebo dva?\\[1 mm]
{\sl Automobilista platil pokutu proto, že jel nedovolenou rychlostí
vesnicí, nikoliv kvůli světlům, protože bylo asi poledne (Měsíc
je v~první čtvrti na obzoru v~poledne).}\\

\noindent
U hospodáře byli tři zájemci o selata. Ve vepřinci byla
právě tři prasátka. Každý si chytil jedno a ještě dvě selata
zbyla. Jak je to možné?\\[1 mm]
{\sl Jeden ze zájemců o sele se jmenoval Každý. Ten chytil sele
a další dva zájemci čekali.}\\

\noindent
V~jedné bitvě byl kdysi zajat nepřátelský generál. Přivedli ho
k~vítězi za veliké slávy, v~přítomnosti celého důstojnického sboru.
Vítězný generál se tvářil velmi blahosklonně a pravil: ,,Nechci být
sám soudcem, ponechám váš život osudu a náhodě. Zítra ráno před
zraky celého shromáždění vám předložím sáček s~dvěma kuličkami,
jednou bílou, druhou černou. Vytáhnete si jednu kuličku. Vytáhnete-li
černou, budete zastřelen, vyjmete-li bílou, propustím vás na svobodu.''
Ale žalářník, který byl vězni nakloněn, mu prozradil, že nemá naději
na záchranu, protože v~sáčku budou obě kuličky černé. Odsouzenec
celou noc přemýšlel, jak se zachránit, aniž by vzbudil nelibost
nepřítele například tím, že by chtěl obě kuličky předem ukázat.
Nakonec přišel na šťastné řešení, které mu zachránilo život.
Jak to provedl?\\[1 mm]
{\sl Když ráno předložili odsouzenci sáček s~kuličkami, vyňal jednu
a aniž by se na ni podíval, spolkl~ji. ,,Podívejte se, jaká zbyla!''
žádal své soudce. Protože zbyla černá kulička, museli mlčky
připustit --- aby záměrný podvod se veřejně neprovalil --- že ta
první kulička nemohla být jiná, než bílá.}\\

\noindent
Letadlo letělo nejprve 100~km směrem k jihu, pak 100~km přímo na
východ a konečně zase 100~km na sever. Pilot se však nesmírně
podivil, že přistál přesně na témž místě, odkud vzlétl.
Které to bylo místo?\\[1 mm]
{\sl To místo je severní pól. Ale nutno tuto úlohu doplnit tím,
že existuje ještě jedno místo, odkud pilot dosáhne téhož
výsledku.}\\

\noindent
Dejme tomu, že letadlo spadne přesně na hranici mezi státem,
odkud pochází letecká společnost havarovaného letadla a
státem, odkud pochází většina cestujících. Který z obou států
dodá rakve pro pozůstalé po obětech?\\[1 mm]
{\sl Stačí pohřbít pozůstatky po obětech. Pozůstalé necháme
pečovat o jejich hroby.}\\

\noindent
Dívám se na čísi podobiznu. Zajímalo by vás, kdo je na ní
zobrazen? Prozradím vám, že nemám sourozence a že otec toho
muže na obrázku je syn mého otce. Kdo je na obrázku?\\[1 mm]
{\sl Úlohu můžeme rozdělit na dvě části. Je zřejmé, že syn mého
otce jsem já sám. Pak celkový výraz lze přepsat do tvaru: Otec
toho muže na obrázku jsem já. Čili na obrázku je můj syn.}\\

\noindent
Nemám sourozence. Syn toho pána na obrázku je syn mého otce.
Kdo je na obrázku?\\[1 mm]
{\sl Můj otec.}\\

\noindent
Jeden paša se chtěl kdysi dobře pobavit a proto uspořádal 
nevšední závod velbloudů. Podmínkou bylo, že velkou odměnu obdrží 
ten, jehož velbloud doběhne k cíli poslední. Nelze se tedy divit, 
že k tak neobvyklému závodu se přihlásili jen dva závodníci. 
Paša i četní diváci očekávali, že jezdci budou zvířata co nejvíce 
zdržovat. Po obvyklém výstřelu, který oznamoval začátek závodu, však
vyrazili oba velbloudi jako diví a poháněni nemilosrdnými ranami holí 
se hnali kolem užaslých diváků k cíli. Víte proč?\\[1 mm]
{\sl Jezdci si před závodem vyměnili velbloudy. Protože v pravidlech
soutěže stálo, že vyhraje ten, jehož velbloud doběhne k cíli
poslední, každý se snažil přijet na velbloudu svého soupeře první.}\\

\noindent
Prvoci v nádobě se dělí každou minutu na dva. Rozdělený
prvok se dělí na dva , z nichž každý má stejný objem jako ten
původní. Nádoba byla plná ve 12.00. Kdy byla z polovice plná?\\[1 mm]
{\sl Za minutu dvanáct.}\\

\noindent
Když můžete koupit 1 za 10 korun, 12 za 20 korun a 150 za
30 korun, co si kupujete?\\[1 mm]
{\sl Kupuji si čísla, kus za 10~Kč.}\\

\noindent
Mám dohromady dvě mince v hodnotě 60 haléřů a jedna z nich 
není padesátihaléř. Jak je to možné?\\[1 mm]
{\sl Padesátihaléř je ta druhá mince.}\\

\noindent
Kolik měsíců má dvacet devět dnů?\\[1 mm]
{\sl 29 dní má 11 měsíců, tedy pokud není přestupný rok, to má
12 měsíců 29 dní.}\\

\noindent
Láhev a zátka stojí korunu deset haléřů. Láhev je o korunu 
dražší než zátka. Kolik stojí zátka?\\[1 mm]
{\sl Pět haléřů.}\\

\noindent
Co je správně: Sedm a pět je třináct nebo sedm a pět jsou 
třináct?\\[1 mm]
{\sl Nic z toho, je to dvanáct.}\\

\noindent
Pán a paní Sedláčkovi mají 6 dcer a každá z nich má bratra. 
Kolik je lidí v této rodině?\\[1 mm]
{\sl Devět.}\\

\noindent
Dva otcové a dva synové zastřelili každý jednu kachnu. Celkem 
zastřelili tři kachny --- jak je to možné?\\[1 mm]
{\sl Děd, otec, syn.}\\

\noindent
Víme že dvanáct tříkorunových známek je do tuctu, ale kolik 
čtyřkorunových známek je v tuctu čtyřkorunových?\\[1 mm]
{\sl Dvanáct.}\\

\noindent
Ve třech krabicích jsou v jedné dvě černé koule, v druhé 
dvě bílé a třetí s bílou a černou. Kolik kuliček musíte vytáhnout 
předem bez toho, abyste viděli jejich barvu, abyste si byli jisti, 
že pak budete schopni určit barvy v jednotlivých krabicích?\\[1 mm]
{\sl Čtyři.}\\

\noindent
Poklad, umístěný 5 metrů od stromu, je hlídán psem u stromu 
uvázaným na 10metrové šňůře. Jak to uděláme, abychom poklad
získali?\\[1 mm]
{\sl Abychom poklad získali, bylo by dle mého názoru třeba obíhat 
se psem v patách okolo stromu, čímž by se 10ti metrová šňůra zkrátila 
natolik, že by nebyl problém poklad získat. (Toto řešení ale 
neplatí v případě, že je smyčka kolem stromu natolik volná, že 
se otáčí a nenamotává nebo také v případě, že rozměr pokladu 
je tak velký, že se pes omotá okolo něj a nebo také v případě, 
ze pes je natolik chytrý, že se sám zpětným pohybem odmotá.)\\[1 mm]
V krajním bodě je možné, že má vzdálenost od pokladu bude 
5m a vzdálenost psa od pokladu 15m (pokud bude na druhé straně 
stromu). Není pak problém pro hbitého člověka poklad vzít a dostat 
se bez újmy z dosahu hlídacího psa. (Ze zadání bohužel nevyplývá, 
zda je možné poklad unést a také není určeno, zda rasa psa je 
běhavá či je to pes gaučový.)\\[1 mm]
Je možné zvolit průkopnickou metodu a k pokladu se prokopat 
zespod. (Bohužel není určeno, zda je půda měkká nebo má skalní 
podklad a také zda se náhodou nejedná o hlídacího psa norníka, 
který by za námi mohl vlézt.)\\[1 mm]
Nabízí se také spousta jiných variant --- krádež ze vzduchu, 
uspání psa, atp. ale je otázkou kterou z nich by si vychytralý
zloděj pokladu zvolil.}\\

\noindent
Jak dokážete, že polovina z dvanácti je sedm?\\[1 mm]
{\sl Pokud napíšeme římskou číslici 12, pak je opravdu její horizontální 
polovinou číslo 7.}\\

\noindent
Jedna mince z devíti je falešná. Máme k dispozici pouze rovnoramenné 
váhy a víme, že falešná mince je lehčí než ostatní. Kolika váženími 
dokážeme falešnou minci odhalit?\\[1 mm]
{\sl 2 vážení: 3+3mince, 1+1 mince}\\

\noindent
Sečeme-li věk muže a ženy, dostaneme součet 91. Jemu je nyní 
dvakrát tolik, kolik bylo jí, když on byl tak stár, jako ona 
nyní. Jaký je věk obou manželů?\\[1 mm]
{\sl Muž má 52 let a žena 39 let.}\\

\noindent
Můžete nakreslit následující obrázek jedním tahem? 100\\[1 mm]
{\sl Ano můžeme. Použijeme k tomu 2 papíry, přičemž na jeho spodní 
část píšeme nutná spojení jednotlivých číslic.}\\

\noindent
Tenisový turnaj začíná se stovkou hráčů. Kolik zápasů je 
nutno odehrát, aby byl znám vítěz?\\[1 mm]
{\sl 99 (50, 25, 12, 6, 3, 2, 1)}\\

\noindent
Máte krabici s dvanácti bílými a dvanácti černými koulemi. 
Kolik jich musíte vytáhnout, abyste si byli jisti, že máte černobílou 
dvojici?\\[1 mm]
{\sl 13}\\

\noindent
Gril umí opéci jednu stranu steaku za 10 minut. Na gril se 
vejdou dva steaky. Jak dlouho trvá opéci tři steaky z obou
stran?\\[1 mm]
{\sl 30}\\

\noindent
Lháři vždy lžou. Pravdomluvní vždy mluví pravdu. Potkáte 
tři muže na cestě a zeptáte se prvního z nich. ,,Jste lhář nebo 
pravdomluvný?'' Nemůžete porozumět jeho odpovědi a tak požádáte 
druhého muže: ,,Co řekl?'' Druhý muž odpoví: ,,On říká, že je 
lhář.'' Třetí muž říká: ,,Druhý muž lže.'' Je třetí muž lhář 
nebo pravdomluvný?\\[1 mm]
{\sl Je pravdomluvný.}\\

\noindent
Žába skáče do 20 metrového kopce a každý den vyskáče o 5 
metrů a v noci o čtyři spadne. Za jak dlouho se dostane na druhou 
stranu kopce?\\[1 mm]
{\sl 16}\\

\noindent
Hřeben střechy směřuje od severu k jihu. Sedne na něj holub 
a snese vejce. Na kterou stranu spadne vejce pravděpodobněji, 
na západní, nebo na východní?\\[1 mm]
{\sl Nespadne nikam, protože holub nesnáší vejce --- snáší je 
jen holubice.}\\

\noindent
V rámci funkční hromadné dopravy ve fiktivním městě autobusy 
přijížděly jeden za druhým. Kolik jich bylo?\\[1 mm]
{\sl 3 (jeden za druhým --- za druhým následuje třetí).}\\

\noindent
Nákladní auto je o 1 cm vyšší než podjezd. Co musíte udělat, 
abyste podjezdem projeli?\\[1 mm]
{\sl Zcela nebo částečně vypustit pneumatiky.}\\

\noindent
2 vlaky, vzdálené 240 km, vyjíždějí ve stejný čas proti sobě 
stejnou rychlostí 60 km/h. Od jednoho z nich v ten samý okamžik 
vyletí moucha rychlostí 75 km/h. Až doletí k protijedoucímu, 
stále se přibližujícímu, vlaku, otočí se a letí nazpět. Takto 
mezi nimi poletuje, dokud se vlaky nesrazí a mouchu nerozmáčknou. 
Kolik kilometrů moucha nalétala?\\[1 mm]
{\sl 150 km (vlakům bude trvat 2 hod. než se srazí, za tyto dvě
hodiny moucha uletí 150 km).}\\

\noindent
V závodě na 400 m porazil Franta Marii v cíli o 80 m. Dohodli 
se, že v opakovaném závodě bude Franta startovat o 100 m za startovní 
čarou. Kdo tento druhý závod vyhraje?\\[1 mm]
{\sl Doběhnou do cíle současně. (V prvním závodě uběhl Franta 400 
a Marie 320m --- Franta uběhne 5/4m ve stejném čase jako uběhne 
Marie 1 m, ve druhém závodě musí běžet Franty 500m --- tuto vzdálenost 
tedy urazí ve stejném čase jako Marie svých 400m.)}\\

\noindent
Z Prahy do Ostravy vyjede vlak. O hodinu později vyjede vlak 
z Ostravy do Prahy. Oba vlaky jedou stejně rychle. Který z vlaků 
bude blíž Ostravě, až se potkají?\\[1 mm]
{\sl Když se potkají, budou od Ostravy oba vlaky ve stejné vzdálenosti 
(pokud však nebudeme brát v úvahu, kterou částí se vlaky potkají 
napřed. Pokud bychom zohlednili to, že se napřed potkají svými 
lokomotivami, tak v tom okamžiku je vlak vyjíždějící z Ostravy 
jako celek k Ostravě blíže).}\\

\noindent
Čtvercový příkop obklopuje hrad. Je deset metrů široký, ale 
vy máte jen dvě devítimetrová prkna. Jak to vyřešíte?\\[1 mm]
{\sl Stlouci (svázat, slepit) prkna k sobě a vytvořit tak jedno
prknodelší než 10m.}\\

\noindent
S pomocí prkna přeskočit na druhou stranu.\\[1 mm]
{\sl Položit jedno prkno tak, aby přečnívalo alespoň 1m přes
příkop, zatížit na jednom konci, na druhý konec položit druhé
prkno a přejít příkop.}\\

\noindent
8 mužů chce každý jeden pokoj v hotelu, ale v hotelu mají 
k dispozici jen 7 pokojů. Recepční říká, že to zařídí a umístí 
dva muže do prvního pokoje, třetího do druhého pokoje, čtvrtého 
do třetího pokoje, pátého do čtvrtého pokoje, šestého do pátého 
pokoje a sedmého do šestého pokoje. Nyní přesune jednoho muže 
z prvního pokoje do sedmého a má vyhráno. Každý muž je sám v 
jednom pokoji. Je tomu tak?\\[1 mm]
{\sl Osmý muž stále čeká v recepci.}\\

\noindent
Muž žije v 19tém poschodí a každý den sjede do práce výtahem. 
Ovšem když se večer sám vrací výtahem, jede do 12tého a pak jde 
pěšky po schodech do 19tého. Proč?\\[1 mm]
{\sl Je trpaslík. (Asi je dost hloupý, jinak by si vzal s sebou 
tyčku, kterou by dosáhl na tlačítka do vyšších pater.)}\\

\noindent
Potkáte dva a nevíte, zda jsou pravdomluvní a nebo lháři. 
Zeptáte se prvního, zda je pravdomluvný. Nerozumíte jeho odpovědi. 
Ale druhý muž odpoví: ,,On řekl ano, ale je velký lhář.'' Kdo 
jsou?\\[1 mm]
{\sl První je lhář a druhý pravdomluvný.}\\

\noindent
Sekretářka napsala čtyři dopisy čtyřem různým lidem a adresy 
na čtyři obálky. Vložila náhodně dopisy do obálek, bez toho, 
aby se na ně dívala. Jaká je šance, že právě tři dostanou správnou
obálku?\\[1 mm]
{\sl Pokud funguje bezchybně pošta, pak je pravděpodobnost nulová.
Kdo by dostal čtvrtou?}\\

\noindent
Jak často potká minutová ručička hodinovou od 12:00 do 24:00 
hodin?\\[1 mm]
{\sl 11 krát}\\

\noindent
Nakreslete čtyři přímé čáry aniž byste zvedli tužku od papíru 
a zároveň protněte všechny body.\\
\begin{picture}(19,19)(-2,-2)
 \multiput(0,0)(5,0){3}{\multiput(0,0)(0,5){3}{\circle*{2}}}
\end{picture}
\begin{picture}(19,19)(-2,-2)
 \multiput(0,0)(5,0){3}{\multiput(0,0)(0,5){3}{\circle*{2}}}
 \put(15,15){\vector(-1,-1){15}}
 \put(0,0){\vector(1,0){15}}
 \put(15,0){\vector(-1,1){15}}
 \put(0,15){\vector(0,-1){10}}
\end{picture}\\

\noindent
Cihla váží 3/4 cihly a 3/4 kilogramu. Kolik kilogramů váží 
cihla?\\[1 mm]
{\sl 3 kg}\\

\noindent
8 koulí vypadá stejně, ale jedna je o trochu těžší než ostatní. 
Jak ji zjistíme na dvě vážení na dvouramenných (vyvažovacích) 
vahách?\\[1 mm]
{\sl 2 vážení: 3+3 koule, 1+1 koule}\\

\noindent
Muž bydlí blízko zastávky, odkud vedou dvě.tramvajová spojení 
do jeho práce. Jedno přes jižní část města a druhé přes severní 
část města. Obě mají stejný 10 minutový interval. Protože to 
trvá úplně stejně, nasedne v náhodnou dobu do první tramvaje, 
kterou vidí. Po roce zjistil, že v průměru na devět cest severním 
městem připadá jedna cesta jižním. Čím to je?\\[1 mm]
{\sl Tramvaj jedoucí severním městem přijíždí minutu po tramvaji 
jedoucí jižním městem.}\\

\noindent
Řidič je známý skrblík a tak ví že na 100 000 kilometrů požil 
každou pneumatiku stejně. Má pět pneumatik včetně rezervní. Jak 
dlouho byla každá z nich v provozu?\\[1 mm]
{\sl 80.000 km}\\

\noindent
Čtvercové okno 5x5 metrů bylo z poloviny plochy zamalováno. 
Přesto zbyla čtvercová nezamalovaná plocha ve dvou různých směrech 
měřící 5 metrů. Jak je to možné?\\[1 mm]
{\sl Směry jsou úhlopříčky čtverce, který je vepsán původnímu 
a je otočen o 45 stupňů.}\\

\noindent
Máte dvoje přesýpací hodiny. Jedny měří 7 minut a druhé 4 
minuty. Jak naměříte 9 minut?\\[1 mm]
{\sl Pustíme oboje hodiny, po 4 minutách 4min hodiny otočíme a 
když se dosypou 7min hodiny, začneme na 4min měřit čas. Po dosypání 
4min hodin je obrátíme a ještě jednou a máme přesně 9 minut (1+4+4).}\\

\noindent
Pan A lže v pondělí, úterý a středu. Pan B lže ve čtvrtek, 
pátek a sobotu. Jeden den pan A říká: ,,Včera byl jeden z mých 
lhavých dnů.'' Pan B říká: ,,Včera byl jeden z mých lhavých dnů 
také.'' Jaký den v týdnu to bylo.\\[1 mm]
{\sl Ve čtvrtek.}\\

\noindent
Potkáte tři lidi a víte, že jsou lháři nebo pravdomluvní. 
První říká: ,,Všichni jsme lháři.'' Druhý říká: ,,Přesně jeden z 
nás je pravdomluvný.'' Kdo jsou?\\[1 mm]
{\sl První je lhář, druhý je pravdomluvný a třetí je lhář.}\\

\noindent
Dělník udělá danou práci za tři dny. Má pomocníka, který 
ovšem dělá za dva dny, co tento dělník za den. Za jak dlouho 
to udělají dohromady?\\[1 mm]
{\sl Za 2 dny.}\\

\noindent
Novák má klece s ptáky a všechny klece jsou obsazeny. Ví, 
že je více klecí, než ptáků v jakékoliv kleci. Musí být pak alespoň 
ve dvou stejný počet ptáků?\\[1 mm]
{\sl Ano. Má-li Novák klec a má-li být obsazená, musí v~ní 
být pták. Má-li ale být více klecí než ptáků v~kterékoli 
z~nich, musí mít dvě klece. Mají-li být obě obsazené různým 
počtem ptáků, potřebuje Novák 1+2 ptáky. To ale už musí mít 3 
klece, atd. Má-li Novák konečný počet klecí, musí být tedy alespoň 
ve dvou stejně ptáků.}\\

\noindent
Skupina letadel je na ostrově. Chtějí obletět zeměkouli, 
ale mají nádrže jen na polovinu cesty. Mohou libovolně přečerpávat 
za letu, ale všechny se musí vrátit na základnu. A jedno z nich 
má obletět zeměkouli.\\[1 mm]
{\sl Letadla stačí tři. První dvě vyletí, až mají 1/4 nádrže prázdné, 
odevzdá druhé letadlo půl nádrže a vrátí se zpět. Pak vyletí 
prvnímu naproti třetí letadlo. Potkají se, až prvnímu dochází 
palivo, třetí odevzdá prvnímu půl nádrže a spolu se vrátí.}\\

\noindent
Jan prodal auto za 40 000. Pak ho koupil za 35 000 a prodal 
ho za 45 000. Kolik na tom vydělal?\\[1 mm]
{\sl 50 000 --- cena auta.}

\noindent
Máme tři skříňky zlatou, stříbrnou a olověnou. Odměna je 
v jedné z nich. Na každé z nich je nápis. Na zlaté stojí ,,Cena 
je v této skříňce.'' Na stříbrné ,,Cena není v této skříňce'' 
a na olověné ,,Cena není ve zlaté skříňce``. Nejvýše jedna z vět 
je pravdivá. Kde je odměna?\\[1 mm]
{\sl Když všechny lžou, zlatá lže ,,tady'' a olověná lže ,,ve zlaté 
ne'', spor. Když zlatá nelže, zlatá nelže ,,tady'' a olověná lže 
,,ve zlaté ne'', spor. Když stříbrná nelže, zlatá lže ,,tady'' 
a olověná lže ,,ve zlaté ne'', spor. Když olověná nelže, stříbrná 
lže ,,tady ne'' a olověná nelže ,,ve zlaté ne'', odměna je tedy 
ve stříbrné.}\\

\noindent
12 koulí jsou podobné, ale jedna je trochu těžší či lehčí. 
Jak zjistíte na tři vážení na rovnoramenných vahách, která to
je?\\[1 mm]
{\sl Zvážím 6+6 koulí. Skupinu s~lehčí koulí rozdělím a 
zvážím 3+3. Ze skupiny s~lehčí koulí odstraním jednu a 
zbylé dvě zvážím. Jsou-li vážené stejné, je lehčí ta nevážená, 
jinak jedna z~nich.}\\

\noindent
Arabský šejk zanechal 17 velbloudů (18tý velbloud náhle zemřel) 
svým třem synům. Odkázal polovinu velbloudů nejstaršímu, třetinu 
mladšímu a devítinu nejmladšímu. Jak to udělat, aby to vyšlo?\\[1 mm]
{\sl Přidám k~velbloudům i mrtvolu. Nejstarší dostane 9, 
mladší 6 a nejmladší 2.}\\

\noindent
Doplňte číslice 0--9 místo písmenek. Každé písmeno reprezentuje 
jinou číslici: KAM + JÍT = DNES.\\[1 mm]
{\sl 789 + 264 = 1053}\\

\noindent
Umístěte čtyři jedničky, tři dvojky a tři trojky do kruhu 
tak, že součet tři po sobě jdoucích čísel bude vždy násobek
tří.\\[1 mm]
{\sl Nemá řešení.}\\

\noindent
Čtyři muži jdou přes řeku. Začínají na jedné straně a mají 
se dostat na druhou za 17 minut. Je noc a mají jen jedno použitelné 
světlo na cestu. Nejvíce dva lidé mohou jít přes most a musí 
jít spolu rychlostí toho pomalejšího. Každý z nich jde jinou 
rychlostí. Muž číslo 1 přejde most za minutu, muž číslo 2 za 
2 minuty, muž číslo 3 za 5 minut a muž číslo 4 za 10 minut. V 
jakém pořadí budou chodit přes most?\\[1 mm]
{\sl 12 -- 1 -- 34 -- 2 -- 12}\\

\noindent
Doplňte na místo písmen číslice 1--6 $A/B+C/D=E/F$.\\[1 mm]
{\sl $1/2 - 4/3 = 5/6$}\\

\noindent
Náhodou jsem potkal dva matematiky T a N. Myslel jsem si 
dvě celá čísla obě větší jak 1 a jejich součet jsem řekl T a 
jejich součin jsem řekl N. Oba vědí podstatu problému, ale neslyšeli 
číslo jež jsem řekl druhému.\\
1. T: Jsem si jist že neznáš mé číslo.\\
2. N: Ještě neznám tvé číslo.\\
3. T: Teď už vím tvé číslo.\\
Jaká čísla jsem si myslel?\\[1 mm]
{\sl 11 a 30}\\

\noindent
Vyjádřete číslo $\pi = 3.1415926535897932$ pomocí číslic 1 až 
9 každé jednou a použitím libovolného počtu operací +-*/ závorek 
a desetinné tečky.\\[1 mm]
{\sl Přesné řešení samozřejmě není možné. Zde jsou nejlepší možnosti:\\
$(7/65.3+.94)*(8/2-1) = 3.14159264$\\
$.8 * (4+5) - 2 * (6/9.1 + 1/.73) \cong 3.14159265392142$\\
$3.2-4/(68+1-5/9.7) \cong 3.14159265392142$}\\

\noindent
Doplňte čísla 1 až 9 tak, aby platilo **/* + **/* + **/* 
= 11 (Nepoužívejte 1 ve spodní část zlomku.) Lze na pravou
stranu dát i čísla 12, 13 či 14?
{\sl $18/9 + 24/6 + 35/7 = 11$\\
12 nejde\\
$19 / 4 + 35 / 7 + 26 / 8 = 13$\\
$38 / 6 + 42 / 7 + 15 / 9 = 14$, nebo 
$42 / 6 + 35 / 7 + 18 / 9 = 14$}\\

\noindent
Sedm mužů jde přes řeku. Začínají na jedné straně a mají 
se dostat na druhou za 25 minut. Je noc a mají jen jedno použitelné 
světlo na cestu. Nejvíce tři lidé mohou jít přes most a musí 
jít spolu rychlostí toho pomalejšího. Každý z nich jde jinou 
rychlostí. Muž číslo 1 přejde most za minutu, muž číslo 2 za 
2 minuty, muž číslo 3 za 6 minut, muž číslo 4 za 7 minut, muž 
číslo 5 za 8 minut , muž číslo 6 za 9 minut a muž číslo 7 za 
10 minut. V jakém pořadí budou chodit přes most?\\[1 mm]
{\sl 12 -- 1 -- 567 -- 2 -- 134 -- 1 -- 12}\\

\noindent
Dá se sestavit pětiúhelník s celočíselnými stranami a celočíselným 
průměrem kružnice opsané?\\[1 mm]
{\sl Strany 5,5,5,8,6}\\

\noindent
Na mapě města je 6 zastávek autobusů A, B, C, E, F a G. Zelené 
autobusy jezdí jen z nebo do těchto stanic, ale vždy jen v uvedeném 
směru. Jejich trasy jsou následující.\\
Z A do B a pak do D.\\
Z E do A nebo také do B.\\
Z C do F.\\
Z A do F.\\
Z D do C.\\
Z A jezdí do D i přímo.\\
Z C do A.\\[1 mm]
Jak můžeme projet všech šest stanic a to tak že každou navštívíme 
právě jednou?\\[1 mm]
{\sl E, B, D, C, A, F}\\
Kam není alespoň jedno spojení (i s přestupy) vyjedeme-li 
ze stanice B.\\[1 mm]
{\sl Do E to nejde.}\\

\noindent
Jeden chodec vyrazil z Přerova do Prahy a víme že tuto vzdálenost 
250 km ujde za 7 dnů. Druhý chodec se vydá proti němu a víme 
že tuto cestu ujde za devět dnů. Za jak dlouho se potkají jdou-li 
proti sobě tou samou trasou.\\[1 mm]
{\sl Rychlost pohybu chodců jsou $250/7$ a $250/9$, a tak jejich společná 
rychlost je $250(1/7+1/9)=250*16/63$ a tak se potkají za $63/16$ 
dne což je o trochu méně než za 4 dny.}\\

\noindent
Dva identické šrouby položíme k sobě, ale hlavičky na opačnou 
stranu. Když horním šroubem otáčíme doprava a spodním také doprava (bráno 
při pohledu od hlavičky šroubu k jeho tělu) budou se šrouby od 
sebe vzdalovat či přibližovat?\\[1 mm]
{\sl Zůstanou ve stejné vzdálenosti.}\\

\noindent
Problém z Rhindského papyru: Myslím si číslo. Když k němu 
přičtu jeho dvě třetiny a odečtu jednu třetinu součtu. Můj výsledek 
je 10. Jaké jsem si myslel číslo?\\[1 mm]
{\sl 9}\\

\noindent
Dorota chce dojít z místa P do místa V a přitom se napít 
z potoka jež teče kousek od ní. Jakou cestou má jít?\\
\begin{picture}(40,20)(0,0)
 \put(0,0){\line(1,0){40}}
 \put(2,3){P}
 \put(30,10){V}
\end{picture}\\[1 mm]
{\sl Nejdříve přeneseme V na druhou stranu podél potoka. Spojíme 
přímkou --- bude určitě nejkratší a pak zakreslíme do původní situace.}\\

\noindent
Zahradní nádrž má čtyři kohoutky. Prvním se napustí za dva 
dny, druhým za tři dny, třetím za čtyři dny a posledním za 6 
hodin. Za jak dlouho se napustí všemi čtyřmi kohoutky otevřenými 
současně?\\[1 mm]
{\sl Předpokládejme že den má dvacet čtyři hodin. Pak se za hodinu 
napustí prvním kohoutkem $1/48$ druhým $1/72$ třetím $1/96$ a
čtvrtým $1/6$. To je dohromady $(6+4+3+48)/288= 61/288$. Nádrž se
napustí plná za $288/61$ hodiny, což je 4 hodiny 43 minuty a
přibližně 17 sekund.}\\

\noindent
Když A a B udělají společně danou práci za 6 dnů. B a C tu 
samou práci za 20 dnů a C spolu s A za 7 a půl dne. Za jak dlouho 
ji udělá každý sám?\\[1 mm]
{\sl A za 8 dnů B za 24 dnů a C za 120 dnů.}\\

\noindent
Dvě čísla jsou taková, že když první dostane od druhého 30 jsou
v poměru 2:1. Když druhé dostane 50 od prvního pak jejich poměr
je 1:3. Jaké to jsou čísla?\\[1 mm]
{\sl 98 a 94}\\

\noindent
Čtyřicet lidí (žen, mužů a dětí) má dohromady čtyřicet zlatých 
mincí. Každý muž vlastní 6 mincí, každá žena 3 mince a každé dítě jednu. 
Kolik mužů, kolik žen a kolik dětí se o mince dělí?\\[1 mm]
{\sl 40 dětí.}\\

\noindent
Každý den v 8 ráno opouští přístav ve Chvaleticích jedna 
loď do Ústí nad Labem a jedna jede z Ústí nad Labem do Chvaletic. 
Cesta z Chvaletic do Ústí nad Labem trvá 7 dnů a 7 nocí. Kolik 
lodí potká daná loď cestou?\\[1 mm]
{\sl 13 lodí na cestě a ovšem i v každém přístavu jednu loď takže
celkem 15.}\\

\noindent
Koupil jsem si spoustu 5 korunových a 7 korunových známek. 
Jakou nejvyšší částku v korunách nejsem schopen pomocí těchto 
známek vylepit?\\[1 mm]
{\sl 23}\\

\noindent
Jaké číslo musí být přičteno současně k 100 a k 20 (to samé 
současně), aby poměr nově vzniklých součtů byl 2:1?\\[1 mm]
{\sl 60}\\

\noindent
Můj syn je pětkrát tak starší než moje dcera. Moje žena je 
pětkrát tak starší než můj syn a já jsem dvakrát starší než moje 
žena. A moje babička, jež slaví dnes 81 narozeniny, je stará 
tak, jako naše roky sečtené dohromady. Jak starý je můj syn?\\[1 mm]
{\sl 5 let}\\

\noindent
Napište 5 stejných číslic tak, aby součet takto vzniklých 
čísel byl 14.\\[1 mm]
{\sl $11+1+1+1=14$}\\

\noindent
Kolikrát se mince otočí o 360 stupňů, když jednou obtočí 
druhou minci?\\[1 mm]
{\sl Dvakrát}\\

\noindent
V této hádance je celkem 10 oznamovacích vět. Kolik z nich 
je pravdivých?\\
Přesně jedna věta je nepravdivá.\\
Přesně dvě věty jsou nepravdivé.\\
Přesně tři věty jsou nepravdivé.\\
Přesně čtyři věty jsou nepravdivé.\\
Přesně pět vět je nepravdivých.\\
Přesně šest vět je nepravdivých.\\
Přesně sedm vět je nepravdivých.\\
Přesně osm vět je nepravdivých.\\
Přesně devět vět je nepravdivých.\\
Přesně deset vět je nepravdivých.\\[1 mm]
{\sl Pradivá věta je jen jedna, a to: ,,Přesně deset vět je
nepravdivých.''První oznamovací věta je totiž: ,,V této hádance
je celkem 10 oznamovacích vět.'' Avšak v hádance je vět 11,
proto prvnívěta a pak dalších devět vět začínajících na ,,Přesně...''
jsou nepravdivé.}\\

\noindent
Součet všech možných trojic čtyř čísel je postupně 22, 24, 
27 a 20. Jaká jsou to čísla?\\[1 mm]
{\sl 9,7,4,11}\\

\noindent
Šest sudů s vínem má objem 15,16,18, 19, 20 a 31 litrů. První 
obchodník si odnesl dva sudy. Druhý den si přišel pro dvakrát 
tak vína. Jaký sud zbyl obchodníkovi?\\[1 mm]
{\sl 20 $(15+18)*2=16+19+31$}\\

\noindent
Kolik rýh je na dlouhohrající desce?\\[1 mm]
{\sl Dvě, na každé straně jedna.}\\

\noindent
Pero, guma a sešit stojí 100 korun. Sešit stojí víc než dvě 
pera, a tři pera stojí více než čtyři gumy. Když tři gumy stojí 
více než sešit, kolik stojí sešit?\\[1 mm]
{\sl 55 (pero 26 a guma 19)}\\

\noindent
Silvestrovská: Máme v řadě sklenici prázdnou, pak tři plné 
a pak dvě prázdné. Kolika sklenic musíme vzít do ruky, aby se 
v řadě střídaly prázdné a plné sklenice?\\[1 mm]
{\sl Jedné. Přeliji obsah třetí do poslední prázdné.}\\

\noindent
Napiš číselnou řadu tak, aby platilo $Y_{N+1} + 2Y_{N} = 0$\\[1 mm]
{\sl 1, -2, 4, -8, 16, ...}\\

\noindent
Veronika a Anežka sbíraly jahody. Celkem nasbíraly 49 hrnečků, 
ale Anežka nasbírala o 5 více než Veronika. Kolik hrnečků každá 
nasbíraly?\\[1 mm]
{\sl Anežka 27, Veronika 22}\\

\noindent
Jaké je nejmenší přirozené číslo (celé, kladné), které při 
dělení 13 dá zbytek 5?\\[1 mm]
{\sl 5}\\

\noindent
Barborka říká, že matka portrétované osoby je dcerou té, 
která zde hovoří. Čí je to portrét?\\[1 mm]
{\sl Vnučky}\\

\noindent
Trpaslík pozval přátele na hostinu a když odcházeli, každý 
si s~každým potřásl rukou. Dohromady si 6krát podali ruce. 
Kolik kamarádů si trpaslík pozval.\\[1 mm]
{\sl Tři kamarádi + trpaslík}\\

\noindent
Pan Kyslík si kupuje čtvercovou zahradu. Pan Dusík si kupuje 
čtvercovou zahradu se stranou $2\times$ delší. Kolikrát víc musí pan 
Dusík zaplatit, je-li za 1m$^2$ stejná cena.\\[1 mm]
{\sl Pan Dusík musí zaplatit $3\times$ tolik.}\\

\noindent
Na lodi je 60 pirátů. Každý 3. má skleněné oko, každý 6. 
dřevěnou nohu a zbytek má hák. Kolik pirátů má hák.\\[1 mm]
{\sl Hák má 30 pirátů.}\\

\noindent
Z~těchto dílů sestavte čtverec\\[1 mm]
\begin{picture}(60,20)(0,0)
 \put(0,0){\line(1,0){10}}
 \put(10,0){\line(0,1){5}}
 \put(10,5){\line(-1,0){5}}
 \put(5,5){\line(0,1){10}}
 \put(5,15){\line(-1,0){5}}
 \put(0,15){\line(0,-1){15}}
 
 \put(15,0){\line(1,0){5}}
 \put(20,0){\line(0,1){10}}
 \put(20,10){\line(1,0){5}}
 \put(25,10){\line(0,1){5}}
 \put(25,15){\line(-1,0){5}}
 \put(20,15){\line(0,1){5}}
 \put(20,20){\line(-1,0){5}}
 \put(15,20){\line(0,-1){20}}

 \put(30,0){\line(1,0){5}}
 \put(35,0){\line(0,1){15}}
 \put(35,15){\line(-1,0){5}}
 \put(30,15){\line(0,-1){15}}
 
 \put(40,0){\line(1,0){5}}
 \put(45,0){\line(0,1){5}}
 \put(45,5){\line(1,0){5}}
 \put(50,5){\line(0,1){5}}
 \put(50,10){\line(-1,0){5}}
 \put(45,10){\line(0,1){5}}
 \put(45,15){\line(-1,0){5}}
 \put(40,15){\line(0,-1){15}}
 
 
 \put(70,0){\line(1,0){20}}
 \put(75,0){\line(0,1){10}}
 \put(75,10){\line(1,0){5}}
 \put(80,10){\line(0,1){5}}
 \put(80,15){\line(-1,0){5}}
 \put(75,15){\line(0,1){5}}
 \put(90,20){\line(-1,0){20}}
 \put(70,20){\line(0,-1){20}}

 \put(90,0){\line(0,1){20}}
 \put(75,5){\line(1,0){15}}
 \put(80,10){\line(1,0){5}}
 \put(85,10){\line(0,1){5}}
 \put(85,15){\line(1,0){5}}
\end{picture}\\

\noindent
Sestavte 7 bodů do 4 řádek po třech.\\[1 mm]
\begin{picture}(30,20)(0,0)
 \multiput(0,0)(0,5){5}{\line(1,0){30}}
 \multiput(15,0)(0,10){3}{\circle*{4}}
 \multiput(5,5)(0,10){2}{\multiput(0,0)(20,0){2}{\circle*{4}}}
\end{picture}\\

\noindent
Které slovo má ve svém kladném i záporném znění tentýž
význam?\\[1 mm]
{\sl Svár a nesvár.}\\

\noindent
Ve kterém českém slově (psáno bez čárek) se vyskytuje 
třikrát po sobě slabika ,,ti``?\\[1 mm]
{\sl Štítiti se.}\\

\noindent
Znáte české slovo, které má na začátku dvě ,,u``?\\[1 mm]
{\sl Uudit.}\\

\noindent
Je to světlé a točí se to. Co je to?\\[1 mm]
{\sl Světlé pivo.}\\

\noindent
Kdy se husa dojí?\\[1 mm]
{\sl Když se sní poslední kousek.}\\

\noindent
Početná skupina výletníků se dostala k~překážce --- k řece, u níž
byla původní lávka stržena. U břehu ležela čtyři 
prkna 4.5 m dlouhá, řeka však byla široká 5 m. Jak lze prken
vytvořit provizorní lávku přes řeku bez použití jakýchkoli dalších
pomůcek?\\[1 mm]
\includegraphics{provlavk}\\

\noindent
Dvě stejné číše jsou naplněny až po čárku, označující 
míru - jedna vínem a druhá limonádou. Vezmeme-li z~číše 
naplněné limonádou plnou lžíci a dáme-li její obsah do číše s~vínem, 
promícháme ji a pak opět plnou lžíci této směsi dáme do limonády, 
bude v~číši s~vínem větší procento limonády než 
procento vína v~číši s limonádou?\\[1 mm]
{\sl V~číši s~vínem bude stejné množství limonády,
jako množství vína v~číši s~limonádou.}\\

\noindent
Kde je na světě místo čtyř světových stran jen jedna 
světová strana?\\[1 mm]
{\sl Na severním pólu je jen jih, neboť kam se podíváme, je to 
vždy jižním směrem. Podobně je na jižním pólu jenom sever.}\\

\noindent
Jak je možno udělat ze dvou tyčí deset, aniž by se zlomily.\\[1 mm]
{\sl Dáme-li dvě tyče přes sebe, vznikne X, římsky 10.}\\

\noindent
Je v~moři více písku nebo ryb?\\[1 mm]
{\sl V~moři je více ryb, písek je pod mořem.}\\

\noindent
,,Závodník snědl tři vejce na lačný žaludek a pak se rozběhl.`` 
Ve větě je chyba. Kde?\\[1 mm]
{\sl Závodník mohl sníst na lačný žaludek jen jedno vejce, při
dalších už nebyl žaludek lačný.}\\

\noindent
V~knihovně jsou zařazeny tři díly jistého šíleného 
románu. Mají 296, 274 a 548 stran. Každý list je silný 0,12 mm 
a jedna deska s~předsádkou 2 mm. Jaká je vzdálenost od 
prvé strany prvého dílu k~poslední straně druhého dílu?\\[1 mm]
{\sl Vzdálenost je 4 mm. Horní deska s~předsádkou prvního 
dílu a spodní deska s~předsádkou druhého dílu.}\\

\noindent
Majitel ranče odkázal majetek svým třem synům, přičemž 
stanovil, že se o 17 koní mají rozdělit tak, že nejstarší syn 
obdrží polovinu, prostřední polovinu a nejmladší devítinu. Synové 
se málem umlátili, protože jim dělení nějak nevycházelo. Sousedi 
na ně zavolali šerifa, který přijel brzy na rychlém ryzáku a 
snadno dědictví rozdělil. Jak to udělal?\\[1 mm]
{\sl Šerif přidal k 17 koním svého ryzáka, takže jich bylo 18. 
Z~nich dal polovinu (9) nejstaršímu, třetinu (6) prostřednímu,
devítinu (2) nejmladšímu synovi a svého koně si vzal zase zpět.}\\

\noindent
V~kolik hodin rozdělují ručičky hodinek ciferník 
na dva půlkruhy, v~nichž je v~každém stejný součet 
hodin (tj. součet čísel od 1 do 12)?\\[1 mm]
{\sl Přibližně v 9:17 a 3:49. Součty činí vždy 39.}\\

\noindent
Kupující platil za klobouk, který stál 50,- stokorunou. 
Protože obchodník neměl nazpět, šel bankovku rozměnit nazpět 
do blízké drogerie a pak návštěvníkovi vrátil 50,- Druhého dne 
se drogista přihlásil, že stokoruna je falešná a obchodník mu 
musel dát jinou. Jakou celkovou škodu obchodník utrpěl?\\[1 mm]
{\sl Přišel celkem o 100,- které musel dát drogistovi místo falešné 
stokoruny. Za klobouk mu zůstalo 50,- od drogisty.}\\

\noindent
Bystrý je-li luštitel náš mladý, neztroskotá v~tomto 
taji. 85 let dohromady dědeček s~vnukem prý mají. O 65 let přitom
děda více má než vnouče jeho. Jak stár je vnuk řešení se hledá,
není to nic nesnadného.\\[1 mm]
{\sl Vnukovi je 10 a dědovi 75.}\\

\noindent
Dá-li Frantík Pepíkovi dvě kuličky, bude mít Pepík dvakrát 
tolik kuliček co Frantík. Dá-li Pepík Frantíkovi dvě kuličky, 
budou mít oba stejný počet. Kolik kuliček má Frantík a kolik 
Pepík?\\[1 mm]
{\sl Frantík má 10 kuliček a Pepík 14.}\\

\noindent
Napište jedním slovem suchá tráva.\\[1 mm]
{\sl Seno.}\\

\noindent
Co je mezi mořem a zemí?\\[1 mm]
{\sl Písmeno ,,a''.}\\

\noindent
Čtyřicet bez deseti je padesát. Jak je to možné?\\[1 mm]
{\sl V~římských číslicích: XL - X = L}\\

\noindent
Dobře se zamyslete a odpovězte. Je dříve dům postaven 
nebo zbourán?\\[1 mm]
{\sl Dům je dříve postaven, neboť bourat jej můžeme až po postavení.}\\

\noindent
,,Dcera se už bude vdávat, ačkoliv jsme slavili pětkrát 
její narozeniny a žádné jsme nevynechali,'' prohlásil tajuplně 
otec. Jak je to možné?\\[1 mm]
{\sl Dcera se narodila 29. února, má tedy narozeniny co 4 roky.}\\

\noindent
Je možné, aby polovina z 12 bylo 7?\\[1 mm]
{\sl V~římských číslicích dostaneme podélným rozpůlením 
čísla XII číslo VII.}\\

\noindent
Kterou rostlinu pozná snadno i slepý?\\[1 mm]
{\sl Kopřivu.}\\

\noindent
Čeho je v~kalamáři nejvíce?\\[1 mm]
{\sl Písmene ,,a''.}\\

\noindent
Dva otcové a dva synové byli na návštěvě. Při svačině 
před ně hostitelka předložila tři čaje a vybídla je, aby pili. 
Všichni najednou se napili. Jak je to možné?\\[1 mm]
{\sl Jsou to dědeček, syn a vnuk (tedy pouze tři osoby).}\\

\noindent
Dvanáct mincí leží jako na číselníku hodin. Vašim úkolem je 
přemístit mince tak, aby od 1 do 6 ležely dvě mince na sobě. 
Mince se přemisťují skokem přes dvě mince ležící vedle.\\[1 mm]
{\sl 12$\rightarrow$3, 7$\rightarrow$4, 10$\rightarrow$6, 8$\rightarrow$1, 
9$\rightarrow$5, 11$\rightarrow$2}\\

\noindent
V kruhu leží 18 zápalek a jako 19 je krabička. Vašim úkolem 
je odkudkoli odpočítat 6 sirek, šestou vzít a počítat dál, až 
zůstane pouze krabička. S krabičkou se pracuje stejně jako se 
sirkami.\\[1 mm]
{\sl Jako první se vezme 11. Zápalka od krabičky.}\\

\noindent
Vyjmenujte 5 po sobě jsoucích dnů, jejichž jméno neobsahuje 
U ani K.\\[1 mm]
{\sl Předevčírem, včera, dnes, zítra, pozítří}\\

\noindent
Jak udělat ze čtyř karetních pětek čtyři čtverky. Některou 
kartou něco zakrýt.\\[1 mm]
{\sl Karty se položí do kříže (kruhu), každá karta zakrývá své 
sousedce jeden znak.}\\

\noindent
Jak položit na stůl tři sirky tak, aby se hlavičkami 
nedotýkaly stolu.\\[1 mm]
{\sl Utvoříme ze sirek trojúhelník, přičemž každá sirka bude za 
hlavičkou opřená o sirku vedlejší.}\\

\noindent
Mají v USA 1. máj?\\[1 mm]
{\sl Ano, vždy po 30. dubnu.}\\

\noindent
Kolik dnů narození zažije průměrně člověk?\\[1 mm]
{\sl Pouze jeden, v porodnici.}\\

\noindent
Některé měsíce mají 31 dnů, kolik jich má 28 dnů?\\[1 mm]
{\sl Všechny, tedy 12.}\\

\noindent
Je těžší kilo železa, nebo kilo peří?\\[1 mm]
{\sl Stejně, pokud vám nespadnou na nohu.}\\

\noindent
Je u nás legální oženit se znova se svou ovdovělou
manželkou?\\[1 mm]
{\sl Ne, protože muž, jehož žena je vdova, je mrtev.}\\

\noindent
Vydělte 30 číslem 1/2 a přičtěte 10. Kolik vám vyšlo?\\[1 mm]
{\sl 70, protože 30 děleno 1/2 je 60.}\\

\noindent
Když na misce leží 3 jablka a dvě odeberete, kolik jablek máte?\\[1 mm]
{\sl Dvě. Ty, co jste odebrali.}\\

\noindent
Doktor vám dal tři pilulky a řekl, že máte brát každou půlhodinu
jednu. Kolik minut uteče od první pilulky než si vezmete poslední?\\[1 mm]
{\sl 60. Začnete první pilulkou, o 30 minut později vezmete druhou
a po 30ti minutách poslední.}\\

\noindent
Farmář měl 17 ovcí a kromě devíti mu všechny pošly. Kolik mu
jich zůstalo?\\[1 mm]
{\sl 9, ostatní přece uhynuly.}\\

\noindent
Kolik zvířat každého pohlaví a druhu vzal Mojžíš na svou Archu?\\[1 mm]
{\sl Žádné, to přece nebyl Mojžíš, ale Noe.}\\

\noindent
Kolik dvoukorun je v tuctu?\\[1 mm]
{\sl 12 dvoukorun v tuctu dvoukorun.}\\

\noindent
Na jezírku rostou lekníny. Každé ráno se na hladině objeví $2\times$
tolik leknínů, než kolik tam bylo minulý den. Jezírko úplně zaroste
za 4 dny. Za jak dlouho bude porostlé jen z $1/4$?\\[1 mm]
{\sl Za 2 dny.}\\

\noindent
Napište pomocí šesti číslic 5 a základních aritmetických operátorů
(+ - * /) číslo 1000?\\[1 mm]
{\sl $(5+5)(5+5)(5+5)=1000$.}\\

\noindent
Doplňte číslice 0--9 místo písmenek. Každé písmeno reprezentuje 
jinou číslici: TO + TO + LÉTO + ALE = LETÍ.\\[1 mm]
{\sl 32 + 32 + 4532 + 041 = 4637}\\

\noindent
Doplňte číslice 0--9 místo písmenek. Každé písmeno reprezentuje 
jinou číslici: HLE + HAD + LEZE + NA = ZEMI.\\[1 mm]
{\sl 473 + 406 + 7383 + 90 = 8352}\\

\noindent
Doplňte číslice 0--9 místo písmenek. Každé písmeno reprezentuje 
jinou číslici: DÁME + SI + DO = NOSU.\\[1 mm]
{\sl 6941 + 52 + 60 = 7053}\\

\noindent
Dvojice robotů stojí před naprosto stejnými dveřmi. Jeden ze strojů na tebe
obrátí pohled a řekne: ,,Pravé dveře jsou ty, které potřebuješ.''\\
,,Neposlouchej ho,'' radí ti druhý, ,,lže.''\\
,,Ale jenom někdy,'' reaguje první.\\
Kterými dveřmi půjdeš? Ze [\ref{gabook1}].\\[1 mm]
{\sl Máš jít pravými.}\\

\noindent
Mám tři zlaté truhlice. V každé z nich jsou tři stříbrné truhlice. V každé
stříbrné jsou dvě olověné truhlice. V každé olověné truhlici jsou tři dřevěné
truhlice, ve kterých už žádné truhlice nejsou. Do každé truhlice jsem vložil
jeden zlaťák a ještě tolik zlaťáků, kolik je v ní celkem všech truhlic. Kolik
mám zlaťáků ve všech truhlicích dohromady? Ze [\ref{gabook2}].\\[1 mm]
{\sl $3((1+3(1+2(1+3)))+3((1+2(1+3))+2((1+3)+3))) = 291$}\\

\noindent
Všichni duchové jsou strašidla a někteří duchové jsou obludy. Některé obludy
jsou lidožrouti, takže všichni lidožrouti jsou strašidla. Je to pravda,
nebo ne? Ze [\ref{gabook2}].\\[1 mm]
{\sl Nelze rozhodnout.}\\

\noindent
Mám dvě dcery. Až bude ta mladší tak stará, jako byla starší, když se mladší
narodila, bude zároveň dvakrát starší než je teď. Jaký je jejich věkový
rozdíl, když jim dohromady je 48 let? Ze [\ref{gabook2}].\\[1 mm]
{\sl Rozdíl činí 24 let.}\\

\noindent
Před půl hodinou bylo tři čtvrtě na osm. Kolik bude za hodinu?\\[1 mm]
{\sl Čtvrt na deset.}\\

\noindent
Ve škole je jedna vyučovací jednotka dlouhá 45 minut. Žáci mají 6 vyučovacích
jednotek denně. Kolik hodin se každý den učí?\\[1 mm]
{\sl Čtyři a půl hodiny.}\\

\noindent
Tramvaj přijíždí na zastávku každých 12 minut. První přijela v 9 hodin a 1
minutu. V kolik hodin přijede třetí tramvaj?\\[1 mm]
{\sl V 9 hodin, 25 minut.}\\

\noindent
Nádraží je od mého domu vzdáleno 6 kilometrů. Na kole jezdím rychlostí
12 km/h. Stihnu vlak, jestliže vyjedu 25 minut před jeho odjezdem?\\[1 mm]
{\sl Nestihnu, cesta na nádraží mi trvá 30 minut.}\\

\noindent
Vlak, který měl 20 minut zpoždění, přijel v 8 hodin a 12 minut. V kolik
hodin měl podle jízdního řádu vlak přijet?\\[1 mm]
{\sl V 7 hodin, 52 minut.}\\

\noindent
Za 6 minut bude tři čtvrtě na pět. Kolik bylo před 9 minutami?\\[1 mm]
{\sl Byly 4 hodiny, 30 minut.}\\

\noindent
Jsi v místnosti s jedinými dveřmi. Stojí zde tři sochy: medvěd, vlk
a ještěr. Klíč od dveří je v tlamě jedné ze soch. Musíš tam sáhnout a
klíč vytáhnout. Pokud sáhneš do tlamy, ve které klíč není, socha ti
ukousne ruku.\\
Pod sochami je 6 celkem nápisů, na každém podstavci dva. Na jednom
podstavci jsou oba nápisy nepravdivé, na dalším jsou oba pravdivé a na
zbývajícím je jeden nápis pravdivý a jeden nepravdivý.\\
U vlka je napsáno: ,,Já klíč nemám,'' a ,,Klíč je v tlamě medvěda.'' U
medvěda stojí: ,,Vlk nemá klíč,'' a ,,Klíč najdeš u ještěra.'' Nápisy pod
ještěrem hlásají: ,,Já klíč nemám,'' a ,,Klíč má vlk.''\\
Do které tlamy sáhneš pro klíč?\\[1 mm]
{\sl Klíč je v tlamě ještěra.}\\

\noindent
Na břehu řeky stojí farmář s ovcí, vlkem a hlávkou zelí. U břehu kotví jedna
malá loďka. Loďka je tak malá, že se do ní s farmářem vleze už buď jen koza,
jen vlk, nebo jen zelí.\\
Farmář chce dopravit všechen svůj majetek na druhou stranu. Protože ale vlci
žerou kozy a kozy mají rády zelí, nesmí zůstat bez farmáře na stejném břehu
vlk s kozou, ani koza se zelím.\\
V jakém pořadí má farmář svůj majetek převézt?\\[1 mm]
{\sl Koza tam -- vlk tam -- koza zpět -- zelí tam -- koza tam.}\\

\noindent
Na břehu řeky stojí 3 misionáři a 3 kanibalové. U břehu kotví dvojmístná
loďka. Všichni se chtějí dostat ve zdraví na druhou stranu. Zůstane-li na
jednom břehu více kabibalů než misionářů, kanbalové misionáře seřerou.\\
V jakém pořadí mají jezdit loďkou?\\[1 mm]
{\sl Dva kanibalové tam -- jeden kanibal zpět -- dva misionáři tam -- 
kabibal s misionářem zpět -- dva kanibalové tam -- jeden kabibal zpět --
dva kanibalové tam.}\\

\noindent
Na břehu řeky stojí pětičlenná rodina (syn, dcera, matka, otec a děda).
Přes řeku vede úzká kláda, přes kterou mohou jít současně nejvýše dvě osoby.
Každá osoba přejde přes kládu za jinou dobu. Synovi to trvá 1 vteřinu, dceři
2, matka přejde kládu za 6 vteřin, otec za 8 a děda za 12. Jdou-li dvě osoby
spolu, jdou rychlostí toho pomalejšího.\\
Protože je noc, musí mít ti, co jdou přes kládu, s sebou petrolejku. V lucerně
je ale palivo jen na 30 vteřin.\\
V jakém pořadí musí členové rodiny chodit přes kládu, aby se všichni stihli
dostat na druhou stranu, než lucerna zhasne?\\[1 mm]
{\sl Syn(1) a dcera(3) tam -- dcera(3) zpět -- otec(8) a děda(12) tam -- 
syn(1) zpět -- syn(1) a matka(6) tam -- syn(1) zpět -- syn(1) a dcera(3)
tam.}\\

\noindent
Aby získal ruku šachové princezny, musí princ na šachovém koni splnit nelehký
úkol. Zvolí si libovolné z políček na obrázku. Z něj pak musí skákat z
políčka na políčko tak, aby navštívil všechna ostatní políčka právě jednou
a posledním skokem se vrátit zpět na výchozí pole.\\
Žádné z políček (kromě prvního) nesmí navštívit dvakrát. V jakém pořadí
musí princ po políčkách skákat? Šachový kůň skáče vždy o dvě políčka dopředu
a jedno políčko doprava nebo doleva.\\[1 mm]
\begin{picture}(110,50)(0,0)
 \put(0,39){\framebox(10,10){}} \put(13,39){\framebox(10,10){}} \put(26,39){\framebox(10,10){}} \put(39,39){\framebox(10,10){}}
 \put(0,26){\framebox(10,10){}} \put(13,26){\framebox(10,10){}} \put(26,26){\framebox(10,10){}} \put(39,26){\framebox(10,10){}}
 \put(0,13){\framebox(10,10){}} \put(13,13){\framebox(10,10){}} \put(26,13){\framebox(10,10){}} \put(39,13){\framebox(10,10){}}
 \put(13,0){\framebox(10,10){}} \put(26,0){\framebox(10,10){}}
 \put(60,39){\framebox(10,10){10}} \put(73,39){\framebox(10,10){13}} \put(86,39){\framebox(10,10){6}} \put(99,39){\framebox(10,10){3}}
 \put(60,26){\framebox(10,10){5}} \put(73,26){\framebox(10,10){2}} \put(86,26){\framebox(10,10){9}} \put(99,26){\framebox(10,10){12}}
 \put(60,13){\framebox(10,10){14}} \put(73,13){\framebox(10,10){11}} \put(86,13){\framebox(10,10){4}} \put(99,13){\framebox(10,10){7}}
 \put(73,0){\framebox(10,10){8}} \put(86,0){\framebox(10,10){1}}
\end{picture}\\

\end{multicols}
\clearpage

% End of file

